\documentclass[runningheads]{llncs}

\usepackage[T1]{fontenc}
\usepackage{graphicx}
\usepackage{amsmath}
\usepackage{amssymb}
\usepackage{stmaryrd}
\usepackage{mathpartir}
\usepackage{booktabs,float}
\usepackage{xspace}
\usepackage{url}

\newcommand{\Lean}{Lean~4\xspace}
\newcommand{\State}{\sigma}
\newcommand{\Loc}{\mathit{Loc}}
\newcommand{\Nat}{\mathbb{N}}
\newcommand{\Int}{\mathbb{Z}}
\newcommand{\aval}[1]{\mathcal{A}\llbracket #1 \rrbracket}
\newcommand{\bval}[1]{\mathcal{B}\llbracket #1 \rrbracket}
\newcommand{\tacost}[1]{\mathcal{T\!A}\llbracket #1 \rrbracket}
\newcommand{\tbcost}[1]{\mathcal{T\!B}\llbracket #1 \rrbracket}
\newcommand{\evals}[4]{\langle #1, #2\rangle \rightarrow^{#3} #4}
\newcommand{\triple}[4]{\{#1\}\;#2\;\{#3 \mid #4\}}
\newcommand{\subst}[3]{#1[#2/#3]}
\newcommand{\upd}[3]{#1[#2 \mapsto #3]}
\newcommand{\cst}[1]{\mathsf{C_{#1}}}
\newcommand{\Valid}{\vDash}
\newcommand{\ValidE}{\vDash_{\!=}}
\newcommand{\BigO}{\mathcal{O}}
\newcommand{\BigTh}{\Theta}
\newcommand{\BigOm}{\Omega}
\newcommand{\code}[1]{\texttt{#1}}

% Filled (black) square for QED, overriding llncs's hollow square.
\renewcommand{\squareforqed}{\ensuremath{\blacksquare}}

\begin{document}

\title{Mechanizing a Resource-Aware Hoare Logic
with Asymptotic Cost Assertions in Lean}
\titlerunning{A Mechanized Resource-Aware Hoare Logic with Asymptotic
Assertions}

\author{Gabriel Ara\'ujo Saldanha\orcidID{0009-0009-1965-7882} \and
Rodrigo Ribeiro\orcidID{0000-0003-0131-5154}}
\authorrunning{G. A. Saldanha and R. Ribeiro}
\institute{Programa de P\'os-Gradua\c{c}\~ao em Ci\^encia da Computa\c{c}\~ao,\\
        Universidade Federal de Ouro Preto, Ouro Preto, MG, Brazil\\
        \email{gabriel.saldanha@aluno.ufop.edu.br},\;
\email{rodrigo.ribeiro@ufop.edu.br}}

\maketitle

\begin{abstract}
Reasoning about the cost of a program is central to computer science,
yet the operation counts behind a complexity bound are usually
established informally, where a single miscounted iteration invalidates
the result.
Resource-aware Hoare logics internalize this reasoning through
cost-annotated triples $\triple{P}{S}{Q}{t}$, but their metatheory
typically rests only on pen-and-paper proofs.
We present a mechanization in
\Lean{} of one such logic, covering its cost-instrumented
semantics, the worst-case upper-bound cost logic with soundness and a
relative-completeness result absent from the source, and its verification
condition generator (VCG).
Our mechanization exposed two errors in the original formulation: an unsound
\code{while} rule whose iteration bound is off by one, and a primitive \code{for}
semantics at odds with the fixed per-iteration cost its own rule charges, both of
which we correct and document.
We further add elementary asymptotic ($\BigO$) cost assertions at the level of
triples, as derivable weakenings of the concrete bounds, and use them to analyse
two examples: a division algorithm in $\BigO(n)$, and the halving loop of binary
search in $\BigO(\log n)$.

\keywords{Resource analysis \and Hoare logic \and Cost semantics \and
Asymptotic complexity \and \Lean{} \and Proof mechanization.}
\end{abstract}

\section{Introduction}\label{sec:intro}

The analysis of algorithms is a cornerstone of computer science.
From the first programming courses onward, students are trained to reason about
\emph{how much} a program consumes, not only \emph{what} it computes: to count
the operations executed as a function of the input size and to summarize that
count as an asymptotic growth rate.
This habit of thought is what lets a computer scientist compare designs, choose
data structures, and predict whether a program will scale, and it is exercised
constantly, in lectures, in textbooks, and at the whiteboard.
Program logics, by contrast, have traditionally certified only \emph{what} a
program computes, leaving the resource side of a specification to the same
informal, pen-and-paper argument that fills an algorithms course.

Making that argument formal is both a foundational and a practical concern.
It is foundational because the operation counts behind an asymptotic bound are
proofs like any other and deserve the scrutiny we already apply to correctness;
it is practical because several domains depend on cost guarantees that must not
be wrong.
Real-time and embedded systems must respect execution-time
budgets~\cite{Wilhelm2008}, and cryptographic implementations must run in time
that does not depend on secret inputs, lest they leak those secrets through a
timing side channel~\cite{Barbosa2021,Almeida2016}.
A range of formal approaches has grown around these needs, from type systems
that infer resource bounds automatically~\cite{HofmannJost2003,Hoffmann2012} to
libraries that make asymptotic $\BigO$ claims first-class objects in a proof
assistant~\cite{Gueneau2018}.
What unites them is the goal of putting the resource usage of a program on the
same rigorous footing as its functional correctness, verified together rather
than argued apart.

Silva's dissertation~\cite{Silva2022} develops one such logic for a standard
imperative \code{while}-language, adapting the cost discipline of the EasyCrypt
cost logic~\cite{Barbosa2021}.
Its judgements are cost-annotated Hoare triples $\triple{P}{S}{Q}{t}$, read: if
$S$ starts in a state satisfying $P$ and terminates, then the final state
satisfies $Q$ and the execution consumes resource $t$.
The dissertation instantiates the discipline to \emph{execution time} as a Hoare
logic of worst-case \emph{upper bounds}, in which $t$ bounds the running time of
$S$.
An annotated program, whose loops carry the oracle information their proofs need,
is discharged by a verification condition generator (VCG) that reduces the
validity of the triple to a set of first-order proof obligations.
The metatheory (soundness of the logic and of the VCG) is established
by pen-and-paper proofs.

Pen-and-paper cost proofs are notoriously error-prone: the cost carried in a
conclusion is a numeric expression that must be tracked through every
rule, and a
single off-by-one in an iteration bound silently invalidates a soundness
argument.
This paper reports on a complete mechanization of Silva's proposed logic in
\Lean{}~\cite{Moura2021} on top of the Mathlib library~\cite{Mathlib2020},
and on an extension of the framework with asymptotic cost assertions.
It is worth noticing that this effort is not a transcription exercise:
reconstructing proofs in a
proof assistant forces every implicit assumption into the open and surfaced two
genuine errors in the source text (Section~\ref{sec:deviations}).

More specifically, the main contributions are:
\begin{itemize}
\item A \emph{full mechanization} of the upper-bound resource logic
(Sections~\ref{sec:background}--\ref{sec:logic}): a
deterministic
cost-instrumented operational semantics; the cost
logic as an inductive
judgement with machine-checked soundness; a
\emph{relative-completeness}
theorem that does not
appear in the original text;
and its VCG with a soundness proof.
\item A \emph{shallow-embedding} design in which assertions
and cost terms are
\Lean{} predicates and functions.
The substitution lemma becomes true by definition,
and completeness is
provable without an expressiveness (G\"odel-style)
argument about an
assertion language.
\item The \emph{identification and correction of two errors}
in the source
text (Section~\ref{sec:deviations}): the \code{while}
rule is unsound because its iteration bound is off by
one relative to the cost it charges, and the primitive
\code{for} semantics is inconsistent with the fixed
per-iteration cost its own rule charges.
Both are documented as explicit, justified deviations.
\item An \emph{extension with asymptotic cost assertions}
(Section~\ref{sec:asymptotic}): an elementary
formalization of $\BigO$,
$\BigOm$, and $\BigTh$ over $\Nat \to \Nat$, and a
validity predicate
that lifts $\BigO$ to the level of triples,
indexed by the size of the
input.
Asymptotic assertions arise as derivable weakenings
of the concrete bounds
already proved, and never replace them.
\item Two worked \emph{examples}
(Section~\ref{sec:examples}): a division
algorithm proved to run in $\BigO(n)$ steps, and the
halving loop of binary search proved to run in
$\BigO(\log n)$ steps, where the logarithmic bound
follows from a variant that advances by one against a
geometrically shrinking value.
\end{itemize}
We describe all results using \LaTeX{} notation; every definition follows its
\Lean{} implementation as closely as possible.
All results depend only on the three standard
axioms of \Lean{}'s logic, which are assumed pervasively by Mathlib and which we
describe in Section~\ref{sec:axioms}.
The complete formalization source code and instructions on how to build it are available on-line~\cite{repo-logic}.

The rest of this paper is organized as follows.
Section~\ref{sec:background} reviews the annotated language, its cost model, and
the cost-instrumented semantics.
Section~\ref{sec:logic} presents the mechanized logic and its metatheory, and
Section~\ref{sec:deviations} discusses the corrected errors.
Section~\ref{sec:asymptotic} introduces asymptotic cost assertions, and
Section~\ref{sec:examples} applies them to the two examples.
Section~\ref{sec:related} surveys related work, and
Section~\ref{sec:conclusion} concludes.

\section{The Annotated Language and its Cost Semantics}\label{sec:background}

\subsection{Syntax and States}

The object language is a standard imperative \code{while}-language over integer
variables and one-dimensional arrays.
Arithmetic expressions $a \in \code{AExp}$ are built from integer constants,
scalar variables $x$, array accesses $x[a]$, the binary operators $+, -, \times,
/, \char`\^{}$, and a bounded summation $\sum_{x=lo}^{hi} a$; boolean
expressions
$b \in \code{BExp}$ are the usual comparisons and connectives.
Commands $S \in \code{Stmt}$ are
\[
S ::= \code{skip} \mid x := a \mid x[a_1] := a_2 \mid S_1; S_2 \mid
\code{if}\;b\;\code{then}\;S_1\;\code{else}\;S_2 \mid
\code{while}\;b\;\code{do}\;S.
\]
A memory \emph{location} is either a scalar $\Loc.\code{var}\;x$ or
an array cell
$\Loc.\code{arr}\;x\;n$, and a \emph{state} is a single total map
$\State : \Loc \to \Int$.
We write $\upd{\State}{\ell}{v}$ for the state that agrees with $\State$ except
at $\ell$, where it takes value $v$.
Unifying scalars and array cells into one location type keeps a state a single
function and makes the substitution lemma below a definitional equality.

\subsection{Cost Model}

The cost of each atomic operation is a parameter of the system, so that the
development is not tied to any particular machine.
A \emph{cost model} $C$ is a record assigning a natural number to each syntactic
construct: $\cst{cst}$ for evaluating a constant, $\cst{var}$ for reading a
variable, $\cst{array}$ for an array access, one constant per arithmetic and
boolean operator, and $\cst{skip}$, $\cst{assignV}$, $\cst{assignA}$ for the
statements.
Every theorem is proved for an arbitrary $C$.
The examples use the \emph{unit model}, in which every constant equals $1$; this
is the cost dictionary under which the concrete numbers in
Section~\ref{sec:examples} are stated.

The evaluation cost of an expression is defined compositionally.
Value evaluation $\aval{a}\State$ and $\bval{b}\State$ are standard; the cost
functions $\tacost{a}\State$ and $\tbcost{b}\State$ sum the atomic costs of the
operators involved.
Only the summation $\sum$ makes a cost depend on the state, since its number of
iterations is a value; every $\sum$-free expression has a constant evaluation
cost, a fact the VCG relies on.

\subsection{Cost-Instrumented Operational Semantics}

Execution is modelled by an inductive big-step relation
$\evals{S}{\State}{t}{\State'}$, an inductive \Lean{} proposition whose numeric
index $t$ is the total resource consumed by running $S$ from $\State$ to
$\State'$.
Selected rules are:
\[
\inferrule*[right=\textsc{Assign}]{\;}
{\evals{x := a}{\State}{\tacost{a}\State + \cst{assignV}}
{\upd{\State}{\Loc.\code{var}\,x}{\aval{a}\State}}}
\]
\[
\inferrule*[right=\textsc{Seq}]
{\evals{S_1}{\State}{t_1}{\State'}
\\ \evals{S_2}{\State'}{t_2}{\State''}}
{\evals{S_1;S_2}{\State}{t_1+t_2}{\State''}}
\]
\[
\inferrule*[right=\textsc{While-T}]
{\bval{b}\State = \code{true} \\ \evals{S}{\State}{t}{\State''} \\
\evals{\code{while}\;b\;\code{do}\;S}{\State''}{t'}{\State'}}
{\evals{\code{while}\;b\;\code{do}\;S}{\State}
{\tbcost{b}\State + t + t'}{\State'}}
\]
The \textsc{Assign} rule charges the evaluation cost of $a$ plus the fixed
assignment cost, and updates the scalar location.
The \textsc{Seq} rule adds the costs of the two sub-executions.
Each \code{while} iteration charges the cost of testing the guard,
$\tbcost{b}\State$, in addition to the cost of the body and of the remaining
iterations; the companion rule \textsc{While-F} charges $\tbcost{b}\State$ alone
when the guard is false.
The relation is deterministic: both the final state and the cost are unique.

\begin{lemma}[\code{Eval.deterministic}]\label{lem:det}
If $\evals{S}{\State}{t_1}{\State_1}$ and
$\evals{S}{\State}{t_2}{\State_2}$,
then $t_1 = t_2$ and $\State_1 = \State_2$.
\end{lemma}
Determinism means that a program run from a given state has a single cost and a
single result, so the cost charged by a triple is a well-defined quantity.

\section{The Mechanized Logic}\label{sec:logic}

\subsection{Assertions, Cost Terms, and Validity}

Assertions and cost terms are given a \emph{shallow} embedding: an
assertion is a
\Lean{} predicate over states and a cost term is a function into the naturals,
\[
\code{Assn} \triangleq \State \to \code{Prop},
\qquad
\code{Cost} \triangleq \State \to \Nat,
\]
with the cost term evaluated in the \emph{initial} state of the triple.
The logical variables of the paper logic (names that fix the initial value of
a program variable across the triple) become ordinary \Lean{} variables
universally quantified around the statement.
Substitution into an assertion is state update,
$\subst{P}{a}{x} \triangleq
\lambda\State.\;P\,(\upd{\State}{x}{\aval{a}\State})$,
so the substitution lemma $\State \vDash \subst{P}{a}{x}
\Leftrightarrow \upd{\State}{x}{\aval{a}\State} \vDash P$ holds by
$\code{Iff.rfl}$: it is true \emph{by definition}, not by an inductive argument.

Validity is \emph{total correctness with a cost bound}:
\[
\Valid\;\triple{P}{S}{Q}{T} \;\triangleq\;
\forall\State.\; P\,\State \Rightarrow
\exists\,\State'\,t.\; \evals{S}{\State}{t}{\State'} \wedge Q\,\State'
\wedge t \le T\,\State,
\]
which quantifies over runs that \emph{terminate} and so expresses \emph{total}
correctness.
The original Definition~10 of the dissertation, read literally, quantifies only
over terminating derivations and is therefore \emph{partial} correctness despite
its name; we keep that reading as a separate predicate for the record and adopt
total correctness, without which the \code{while} rule cannot be complete
(Section~\ref{sec:deviations}).

\subsection{The Hoare Kernel and its Metatheory}

The upper-bound logic is an inductive judgement, $\code{Hoare}\;C\;P\;S\;Q\;T$,
whose five arguments are the cost model $C$ (Section~\ref{sec:background}); the
precondition $P$ and the postcondition $Q$, both assertions $\State \to
\code{Prop}$; the command $S$; and the cost term $T : \State \to \Nat$ that bounds
the resource $S$ consumes.
We display this judgement in the customary triple notation, the two denoting the
same object:
\[
\code{Hoare}\;C\;P\;S\;Q\;T
\;\;\equiv\;\;
{\vdash_{C}}\;\triple{P}{S}{Q}{T},
\]
read: under cost model $C$, every terminating run of $S$ from a state satisfying
$P$ ends in a state satisfying $Q$ and costs at most $T$, the bound being
evaluated at the initial state.
Since the cost model $C$ is fixed throughout, we omit it from the notation from
now on, writing simply $\triple{P}{S}{Q}{T}$.
Its atomic and structural rules charge the cost
model directly:
\[
\inferrule*[right=\textsc{Skip}]{\;}{\triple{P}{\code{skip}}{P}{\cst{skip}}}
\qquad
\inferrule*[right=\textsc{Assign}]{\;}
{\triple{\subst{Q}{a}{x}}{x := a}{Q}{\tacost{a} + \cst{assignV}}}
\]
\[
\inferrule*[right=\textsc{If}]
{\triple{P \wedge b}{S_1}{Q}{T_1} \\ \triple{P \wedge \neg
b}{S_2}{Q}{T_2}}
{\triple{P}{\code{if}\;b\;\code{then}\;S_1\;\code{else}\;S_2}{Q}
{\max(T_1, T_2) + \tbcost{b}}}
\]
The interesting rule is \code{while}.
Rather than adopt the dissertation's oracle rule directly, the kernel uses an
invariant $I$, a variant $f : \State \to \Nat$ that strictly decreases on every
iteration, and a \emph{remaining-cost} function $E$ that upper-bounds the cost
still to be paid:
\[
{\small
\inferrule*[right=\textsc{While}]
{ \forall k\,n\,c.\;
        \triple{I \wedge b \wedge f = k \wedge E = n
        \wedge \tbcost{b} + T = c}
        {S}
        {I \wedge f < k \wedge E + c \le n}{T}
        \\\\
        \forall\State.\; I\,\State \wedge
        \bval{b}\State = \code{false}
\Rightarrow \tbcost{b}\State \le E\,\State }
{\triple{I}{\code{while}\;b\;\code{do}\;S}{I \wedge \neg b}{E}}
}
\]
The body premise threads three auxiliary values through the triple: the current
variant value $k$, the current remaining budget $n$, and the cost $c$ of one
guard test plus the body's charged cost $T$.
Its postcondition demands that the variant strictly decreases ($f < k$) and that
the new remaining cost plus what was just paid stays within budget
($E + c \le n$).
The exit premise requires that, once the guard is false, the remaining cost
covers the final guard test.
This formulation is sound and, unlike the paper rule, drives the completeness
proof; the dissertation's oracle rule (its Fig.~4.8) is recovered as a
\emph{derived} rule.

The consequence rule conditions cost weakening on the precondition, which is
sound and stronger than the textbook rule:
$\triple{P'}{S}{Q'}{T'}$ with $P \Rightarrow P'$, $Q' \Rightarrow Q$, and
$\forall\State.\,P\,\State \Rightarrow T'\,\State \le T\,\State$ yields
$\triple{P}{S}{Q}{T}$.
A rule for existentially quantified preconditions,
$\bigl(\forall i.\;\triple{P\,i}{S}{Q}{T}\bigr) \Rightarrow
\triple{\exists i.\,P\,i}{S}{Q}{T}$, is added to support completeness.

\begin{theorem}[\code{Hoare.sound}, Theorem~3
of~\cite{Silva2022}]\label{thm:sound}
If $\vdash \triple{P}{S}{Q}{T}$ then $\Valid\;\triple{P}{S}{Q}{T}$.
\end{theorem}
The proof is by induction on the derivation.
Each atomic rule matches the corresponding \code{Eval} constructor; the
\code{while} case uses the variant $f$ for a well-founded argument on the number
of iterations and the budget invariant $E + c \le n$ to bound the accumulated
cost.

\begin{theorem}[\code{Hoare.complete}]\label{thm:complete}
If $\Valid\;\triple{P}{S}{Q}{T}$ then $\vdash \triple{P}{S}{Q}{T}$.
\end{theorem}
This converse is \emph{not} proved in the dissertation, and it is worth being
precise about \emph{why} our proof carries no expressiveness hypothesis.
Cook's classical relative-completeness theorem~\cite{Cook1978} holds only
\emph{relative} to the \emph{expressiveness} of the assertion language: that for
every command and postcondition the language can name the corresponding strongest
postcondition.
In a deep embedding, where assertions are formulas of a fixed syntactic language,
this is a genuine theorem: the strongest postcondition of a loop quantifies
over the unboundedly many intermediate states of its runs, and capturing that in a
finite formula needs a G\"odel-style arithmetization.
Our shallow embedding sidesteps it entirely: since an assertion \emph{is} an
arbitrary predicate $\State \to \code{Prop}$, the assertion ``language'' is the
full ambient logic of \Lean{}, so the strongest postcondition of a command
(literally $\lambda\State'.\;\exists\,\State\,t.\;P\,\State \wedge
\evals{S}{\State}{t}{\State'}$) and the exact remaining cost of a loop are
\emph{already} predicates, with nothing to encode.
Expressiveness therefore holds by construction, which is precisely what turns
Cook's \emph{relative} converse into an unconditional one here.
The proof uses this directly: for each initial state $\State_0$ satisfying $P$,
validity (with determinism) yields the unique run of $S$, from which we build a
triple with the singleton precondition $\State = \State_0$; the
existential-precondition rule then reassembles this pointwise family into
$\vdash \triple{P}{S}{Q}{T}$.

\subsection{The Verification Condition Generator}

Proving triples by hand is tedious; the dissertation therefore provides a
verification condition generator (VCG) that operates on \emph{annotated}
programs.
The obstacle to full automation is the loop: soundly bounding the cost of a
\code{while} requires facts about it (an invariant, a termination measure, and
a count of how many times the body runs) that are not, in general, computable
from the program text.
Following the dissertation~\cite{Silva2022}, we treat this information as supplied
from outside, by an \emph{oracle} (the programmer, or a separate inference
procedure); the VCG
takes it on trust and reduces the triple to first-order conditions that
\emph{check} the oracle's guesses against the semantics.
Concretely, each \code{while} loop is annotated with a tuple
$(I, f, N, t, \cst{b})$:
\begin{itemize}
\item the \emph{invariant} $I : \State \to \code{Prop}$, a property preserved by
        every iteration and strong enough to entail the postcondition once the guard
        fails;
\item the \emph{variant} $f : \State \to \Nat$, a progress measure that strictly
        advances toward the bound $N$ on each iteration, which is what forces the loop
        to terminate;
\item the \emph{iteration bound} $N : \Nat$: while the guard holds the variant stays
        strictly below $N$, so the body runs at most $N$ times;
\item the \emph{per-iteration cost} $t : \Nat \to \Nat$, the cost charged to the body
        on iteration $k$; making it a function of the index $k$ lets the accounting sum
        costs that differ from one iteration to the next (both examples here use a
        constant $t$, but a body that does more work on later iterations is equally
        expressible);
\item the \emph{guard cost} $\cst{b} : \Nat$, the constant cost of evaluating the
        loop condition once, incurred $N+1$ times: once before each of the $N$
        iterations and once for the final failing test.
\end{itemize}
None of these is trusted blindly: the side conditions below discharge every one of
them against the actual semantics, so a wrong annotation merely makes some
verification condition unprovable instead of yielding an unsound triple.
By recursion on the annotated program $c$, the VCG computes a
weakest-precondition assertion paired with a cost bound, written
$\code{wpc}(c, Q)$, and a finite set of side conditions $\code{VC}(c, Q)$.
Since $\code{wpc}(c, Q)$ is a pair, we write $\code{wpc}(c, Q)_1$ for its first
component, the precondition assertion, and $\code{wpc}(c, Q)_2$ for its second, the
cost bound.
On atomic and compositional commands both are the expected ones; for instance
$\code{wpc}(x := a,\, Q) = (\subst{Q}{a}{x},\ \tacost{a} + \cst{assignV})$, and a
sequence adds the two sub-costs.
The loop case carries the cost accounting:
\[
\code{wpc}(\code{while}\;b\;[I,f,N,t,\cst{b}]\;\code{do}\;c,\ Q)
\;=\; \Bigl(I,\ \ \lambda\State.\ \textstyle\sum_{i < N} t(i) + (N+1)\cdot\cst{b}\Bigr),
\]
charging the $N$ per-iteration body costs plus the $N+1$ guard evaluations.
Its side conditions are
\[
{\small
\begin{array}{rl}
\code{VC}(\code{while}\ldots,\, Q) =
& \bigl(\forall k\,\State.\; I \wedge b \wedge f = k \Rightarrow
  \code{wpc}(c,\, I \wedge k < f)_1\,\State \\[2pt]
& \qquad\qquad {}\wedge\ \code{wpc}(c,\, I \wedge k < f)_2\,\State \le t(k)\bigr) \\[2pt]
& {}\wedge\ \bigl(\forall\State.\; I \wedge \neg b \Rightarrow Q\,\State\bigr)
  \ \wedge\ \bigl(\forall\State.\; I \wedge b \Rightarrow f\,\State < N\bigr) \\[2pt]
& {}\wedge\ \bigl(\forall\State.\; I \Rightarrow \tbcost{b}\State = \cst{b}\bigr)
  \ \wedge\ \code{VC}(c,\, I \wedge k < f).
\end{array}}
\]
Reading top to bottom: the body preserves the invariant and \emph{increases} the
variant $f$ (a progress measure climbing toward the bound $N$, dual to the kernel
rule's decreasing variant) while costing at most $t(k)$; a false guard under the
invariant yields the postcondition; the variant stays below $N$; the guard cost is
the announced constant; and the body's own conditions hold.

Soundness of the VCG requires the annotated command to be \emph{$\sum$-free}: the
summation $\sum$ must not occur in any \emph{executable} expression of the program.
An arithmetic expression is $\sum$-free when it is built without the $\sum$
constructor, that is, from constants, variables, array accesses, and the arithmetic
operators applied to $\sum$-free sub-expressions; a boolean expression is
$\sum$-free when its comparisons relate $\sum$-free arithmetic expressions.
Writing $\code{sf}(\cdot)$ for this predicate on arithmetic expressions, boolean
expressions, and commands, $\sum$-freeness of a command $c$ is defined structurally
by
\[
\inferrule*[right=\textsc{SF-Skip}]{\;}{\code{sf}(\code{skip})}
\qquad
\inferrule*[right=\textsc{SF-Assign}]{\code{sf}(a)}{\code{sf}(x := a)}
\qquad
\inferrule*[right=\textsc{SF-ArrAssign}]
{\code{sf}(a_1) \\ \code{sf}(a_2)}{\code{sf}(x[a_1] := a_2)}
\]
\[
\inferrule*[right=\textsc{SF-Seq}]
{\code{sf}(c_1) \\ \code{sf}(c_2)}{\code{sf}(c_1; c_2)}
\qquad
\inferrule*[right=\textsc{SF-If}]
{\code{sf}(b) \\ \code{sf}(c_1) \\ \code{sf}(c_2)}
{\code{sf}(\code{if}\;b\;\code{then}\;c_1\;\code{else}\;c_2)}
\]
\[
\inferrule*[right=\textsc{SF-While}]
{\code{sf}(b) \\ \code{sf}(c)}
{\code{sf}(\code{while}\;b\;[I,f,N,t,\cst{b}]\;\code{do}\;c)}
\]
Only the executable parts are constrained: the loop annotation
$(I, f, N, t, \cst{b})$ is left unrestricted, so $\sum$ may still appear freely in
assertions and in the annotations, which is exactly where it is needed.

\begin{theorem}[\code{vcg\_sound}, Theorem~1 of~\cite{Silva2022}]\label{thm:vcg}
If $c$ is $\sum$-free and every condition in $\code{VC}(c, Q)$ holds, then the
triple $\triple{\code{wpc}(c,Q)_1}{S}{Q}{\code{wpc}(c,Q)_2}$, where $S$ is the
program underlying $c$, is derivable and hence valid.
\end{theorem}
The $\sum$-freeness hypothesis makes every expression-evaluation cost constant,
which the compositional \code{seq} rule needs (Section~\ref{sec:background}); it
holds of every real program, since $\sum$ occurs only in assertions and
annotations.

\subsection{Trusted Base and Axioms}\label{sec:axioms}

The development rests on exactly three axioms, all part of \Lean{}'s logic and
inherited by anything built on Mathlib, none of them an assumption we chose to
add: \emph{propositional extensionality} (\code{propext}), \emph{quotient
soundness} (\code{Quot.sound}), and \emph{choice} (\code{Classical.choice}).
Here we make precise \emph{where} each one is actually needed.
Checking that a proof uses these and nothing more (via \code{\#print axioms}) is
the standard way to certify that a \Lean{} development is free of hidden gaps.

Quotient soundness (\code{Quot.sound}) enters through two channels.
First, the cost that the VCG computes for a \code{while} loop is a
\emph{finite sum} $\sum_{i < N} t(i)$ over $\code{Finset.range}\;N$, and
\code{Finset} is a quotient of \code{Multiset}, itself a quotient of
\code{List}; so every manipulation of these sums, such as the
$\code{Finset.sum\_const}$ and $\code{Finset.card\_range}$ steps that close the
division example, rests on \code{Quot.sound}.
Second, function extensionality (\code{funext}), which we use whenever two cost
terms of type $\State \to \Nat$ are proved equal, is itself derived from
\code{Quot.sound} in \Lean{}'s kernel.

Propositional extensionality (\code{propext}) is unavoidable because assertions
are shallow, $\code{Assn} = \State \to \code{Prop}$.
Every rewriting step that replaces a proposition by a logically equivalent one,
which is what the pervasive \code{simp} and \code{rw} calls in the rule and
example proofs do, turns an $\leftrightarrow$ into an $=$ and therefore invokes
\code{propext}; it also underlies \code{funext} at \code{Prop}-valued functions
and many of the reused Mathlib order and arithmetic lemmas.

Choice (\code{Classical.choice}) is needed most visibly by the completeness
proof of Theorem~\ref{thm:complete}.
Turning an arbitrary valid triple into a derivation proceeds by case analysis on
whether a given state satisfies the precondition (a \code{by\_cases}
on $P\,\State$
in the \code{of\_pointwise} lemma), and since $P\,\State$ is an arbitrary,
undecidable proposition this excluded-middle step exists only classically.
The same classical footing is what the Mathlib automation (\code{omega} and the
order and arithmetic lemmas over $\Nat$ and $\Int$) already assumes.

By contrast, the lower layers, such as determinism of the semantics
and the soundness theorems, are constructive in
spirit; the axioms surface in them only because they, too, import Mathlib.
The point of reporting exactly these three is thus not that every proof needs
all of them, but that the development introduces no axiom beyond Mathlib's own.

\section{Two Corrections to the Original Formulation}\label{sec:deviations}

Mechanization turned two informal steps of the dissertation~\cite{Silva2022} into failed proof
obligations, each revealing a genuine error. Both had to be fixed for the
corresponding soundness results to hold.

\paragraph{The \code{while} rule is off by one.}
The oracle rule for \code{while} (Fig.~4.8) reads, in our notation,
\[
{\small
\inferrule*[right=\textsc{While-4.8}]
{ \forall k.\ \triple{I \wedge b \wedge f = k}{S}{I}{t(k)}
  \\ I \wedge b \Rightarrow f \le N }
{ \triple{I}{\code{while}\;b\;\code{do}\;S}{I \wedge \neg b}
  {\textstyle\sum_{i=0}^{N-1} t(i) + (N+1)\cdot\tbcost{b}} }
}
\]
where $f$ is the variant, $N$ its announced bound, and $t(k)$ the cost of the body
on the iteration at which $f = k$.
The side condition $f \le N$ admits entry states with $f = 0, 1, \ldots, N$, that
is, up to $N+1$ iterations, whereas the charged cost sums only $N$ body executions
($i$ ranging from $0$ to $N-1$) alongside $N+1$ guard tests.
The honest cost of $N+1$ iterations is $\sum_{i=0}^{N} t(i) + (N+2)\cdot\tbcost{b}$,
one body cost and one guard test more; the rule as written therefore charges too
little and is \emph{unsound}. Indeed the dissertation's own soundness proof
(p.~26) derives this larger bound, not the one the rule concludes.
The fix is to make the bound \emph{strict}, replacing $f \le N$ by
$I \wedge b \Rightarrow f < N$, so that at most $N$ iterations occur and the
announced cost is exactly right.
Our \code{while} rule (Section~\ref{sec:logic}) bakes this in through the
strict decrease $f < k$, and the corrected Fig.~4.8 rule is derived from it.

\paragraph{The primitive \code{for} contradicts its own cost rule.}
The exact-cost fragment of Chapter~6, which the artifact mechanizes, replaces
\code{while} by a bounded \code{for} loop whose big-step semantics (Fig.~6.1)
advances by incrementing the \emph{lower bound}:
\[
{\small
\inferrule*[right=\textsc{For-6.1}]
{ \aval{a}\State < \aval{b}\State
  \\ \evals{S}{\upd{\State}{i}{\aval{a}\State}}{t}{\State'}
  \\ \evals{\code{for}\;i = a{+}1\;\code{to}\;b\;\code{do}\;S}{\State'}{t'}{\State''} }
{ \evals{\code{for}\;i = a\;\code{to}\;b\;\code{do}\;S}{\State}
  {\tacost{a} + t + t'}{\State''} }
}
\]
After $k$ unfoldings the lower bound is the expression $a + 1 + \cdots + 1$ with
$k$ summands, so the cost $\tacost{a + 1 + \cdots + 1} = \tacost{a} +
k\cdot(\cst{cst} + \cst{+})$ it charges to reach the counter \emph{grows} with $k$.
The exact-cost rule for this loop (Fig.~6.2), however, charges a \emph{fixed}
per-iteration expression cost $\tacost{a}$, so the chapter's central exact-cost
result does not hold of the very semantics it is stated over.
We resolve this by making \code{for} syntactic sugar, with the counter kept in the
state,
\[
\code{for}\;i = a\;\code{to}\;b\;\code{do}\;S \;\triangleq\;
i := a;\ \code{while}\;i < b\;\code{do}\;(S;\ i := i+1),
\]
so that each iteration reads $i$ at the constant cost $\cst{var}$ and the
exact-cost accounting becomes faithful. A side effect is that the concrete
constants of some examples shift; each mechanized example states the bound it
actually proves.

\section{Asymptotic Cost Assertions}\label{sec:asymptotic}

The dissertation deliberately avoids asymptotic notation: its stated advantage
over prior order-of-magnitude systems is that it delivers \emph{concrete} costs,
which are strictly more informative than a $\BigO$ class.
Asymptotic assertions are nonetheless the vocabulary in which algorithmic
complexity is usually communicated, and they are often all one needs.
We therefore add them as an extension, taking care that they arise as
\emph{derivable weakenings} of the concrete bounds and never replace them.

\subsection{An Elementary Formalization of $\BigO$, $\BigOm$, $\BigTh$}

We adopt the textbook definitions of complexity classes over functions
$f, g : \Nat \to \Nat$, kept entirely within $\Nat$ so that the example proofs
stay elementary and avoid the normed-space and filter machinery of the general
Mathlib \code{IsBigO}:
\[
f \in \BigO(g) \;\triangleq\;
\exists\,c\,n_0.\; 0 < c \wedge \forall n \ge n_0.\; f\,n \le
c \cdot g\,n,
\]
with $f \in \BigOm(g) \triangleq g \in \BigO(f)$ and
$f \in \BigTh(g) \triangleq f \in \BigO(g) \wedge f \in \BigOm(g)$.
Two facts suffice for the examples: any $\lambda n.\,a\cdot n + b$ with $a > 0$ is
in $\BigO(n)$ (witness $c = a+b$, $n_0 = 1$), and any
$\lambda n.\,a\cdot\log_2 n + b$ with $a > 0$ is in $\BigO(\log n)$ (threshold
$n_0 = 2$, where $\log_2 n \ge 1$).

To make the mechanized text read as in a textbook, each class is additionally
realized as the \emph{set} of functions satisfying the predicate above (\Lean{}
\code{bigO} and its $\BigOm$/$\BigTh$ analogues), equipped with scoped notations
$\BigO(g)$, $\BigOm(g)$, $\BigTh(g)$.
An asymptotic assertion is then literally a membership $f \in
\BigO(g)$, which the
lemma \code{mem\_bigO} unfolds to the predicate above by reflexivity, so the
notation carries no proof overhead.
All the statements below use this membership form.

\subsection{Lifting Asymptotics to Triples}

An assertion that a \emph{number} is $\BigO(g)$ says little about a program.
To state that a \emph{program runs in} $\BigO(g)$ time, we index a family of
pre- and post-conditions by the size $n \in \Nat$ of the input and require a
single cost bound $T : \Nat \to \Nat$, itself in $\BigO(g)$, that is valid at
every size:
\[
\begin{array}{l}
\code{BigOValid}\;C\;P\;S\;Q\;g \;\triangleq\;
\exists\,T : \Nat \to \Nat.\;\\
\qquad T \in \BigO(g) \;\wedge\;
\forall n.\; \Valid\;\triple{P\,n}{S}{Q\,n}{\lambda\_. \,T\,n}.
\end{array}
\]
Reusing the triple notation with the complexity class placed where a concrete
cost term would sit, we abbreviate this predicate as
\[
\Valid\;\triple{P}{S}{Q}{\BigO(g)}
\qquad\text{for}\qquad
\code{BigOValid}\;C\;P\;S\;Q\;g,
\]
where $P, Q : \Nat \to \code{Assn}$ are the size-indexed families and, as in
Section~\ref{sec:logic}, the fixed cost model $C$ is left implicit.
A program is thus certified to run in $\BigO(g)$ time exactly when the cost bound
proved for it, read as a function of the input size $n$, lies in $\BigO(g)$: a
derivable weakening of the concrete bound, obtained without leaving the
upper-bound logic.

\section{Examples}\label{sec:examples}

\subsection{Division: an $\BigO(n)$ Upper Bound}

The first example is the classic Euclidean-style division algorithm, which
computes the quotient $q$ and remainder $r$ of $n$ by $m$ by repeated
subtraction:
\[
\begin{array}{l}
\{\,r = n \wedge q = 0 \wedge m > 0 \wedge n \ge 0\,\}\\
\quad \code{while}\;m \le r\;\code{do}\;(r := r -
m;\; q := q + 1)\\
\{\,n = r + m\cdot q \wedge r < m \mid 20n + 5\,\}
\end{array}
\]
Running the VCG on the annotated program, with invariant
$I \equiv n = q\cdot m + r \wedge m > 0 \wedge r \ge 0$, variant $f
\equiv n - r$,
and iteration bound $N \equiv n$, produces all discharged side conditions and a
computed cost of $13n + 3$ under the unit model, which is weakened to the
announced $20n + 5$.

\begin{theorem}[\code{division\_valid}]
For every initial dividend $n$, the division triple above is valid:
$\Valid\;\triple{P}{S}{Q}{\lambda\_. \,20n + 5}$.
\end{theorem}
The linear bound $20n + 5$ is in $\BigO(n)$, so the concrete result lifts
directly.

\begin{corollary}[\code{division\_bigO}]
The division algorithm runs in $\BigO(n)$ steps:
\[ \Valid\;\triple{P}{S}{Q}{\BigO(\lambda n.\,n)}. \]
\end{corollary}

\subsection{Binary Search: an $\BigO(\log n)$ Bound}

The second example is the logarithmic loop at the heart of binary search:
repeated halving of a value until it reaches one.
\[
\begin{array}{l}
\{\,m = n \wedge n \ge 0\,\}\\
\quad \code{while}\;m > 1\;\code{do}\;m := m / 2\\
\{\,m \le 1 \mid 7\log_2 n + 5\,\}
\end{array}
\]
Each iteration halves $m$, so the loop runs $\lfloor\log_2 n\rfloor$
times: this is
the very iteration profile that makes binary search logarithmic,
isolated from the
array lookup.
What makes the mechanization interesting is the \emph{variant}.
Running the VCG with invariant $I \equiv 0 \le m \le n$, variant
$f \equiv \log_2 n - \log_2 m$, and iteration bound $N \equiv \log_2
n$, the crucial
side condition is that $f$ strictly increases on every iteration.
It does, and by exactly one, because dividing by the base drops the integer
logarithm by one, $\log_2(m / 2) = \log_2 m - 1$ (the Mathlib lemma
$\code{Nat.log\_div\_base}$).
The variant therefore reaches its bound $\log_2 n$ in $\log_2 n$
steps even though
it is measured against a value $m$ that shrinks \emph{geometrically}, so the
$\sum_{i < N}$ that the VCG charges runs over only $\log_2 n$ iterations.
The computed cost under the unit model is $7\log_2 n + 3$, weakened
to the announced
$7\log_2 n + 5$.

\begin{theorem}[\code{halving\_valid}]
For every initial value $n$, the loop terminates with $m \le 1$ and costs at
most $7\log_2 n + 5$:
$\Valid\;\triple{P}{S}{Q}{\lambda\_. \,7\log_2 n + 5}$.
\end{theorem}
Indexing the input by its size $n \in \Nat$, the bound $7\log_2 n + 5$ is in
$\BigO(\log n)$ (the elementary lemma $\code{isBigO\_log}$), so the
concrete result
lifts directly.

\begin{corollary}[\code{halving\_bigO}]
Repeated halving runs in $\BigO(\log n)$ steps:
\[ \Valid\;\triple{P}{S}{Q}{\BigO(\lambda n.\,\log_2 n)}. \]
\end{corollary}
Where the division bound comes from a variant that decreases by one against a
value shrinking arithmetically, the logarithmic iteration count here
comes from a
variant that advances by one against a geometric decrease; the asymptotic layer
records the difference as an $\BigO(\log n)$ class.

\section{Related Work}\label{sec:related}

\paragraph{Program logics for cost.}
Attaching a resource to a Hoare triple $\triple{P}{S}{Q}{t}$ extends the partial
correctness of Hoare~\cite{Hoare1969} with a quantitative component, an idea that
goes back to the earliest work on reasoning about execution time.
Nielson~\cite{Nielson1987} gives a Hoare-like system that certifies the
\emph{order of magnitude} of the worst-case time; it is compositional and clean,
but, as the dissertation we mechanize stresses, an order-of-magnitude judgement
cannot express the concrete, constant-factor bounds that a cost dictionary yields
and that our examples state.
Carbonneaux et al.~\cite{CarbonneauxHoffmannShao2015} pair a quantitative Hoare
logic with an abstract-interpretation front end that \emph{infers} compositional
polynomial resource bounds for imperative code, and certify the whole pipeline in
Rocq.
Haslbeck and Nipkow~\cite{HaslbeckNipkow2018} carry out a comparative
meta-theoretic study of three styles of time-bound Hoare logic (a Nielson-style
system, a big-step system, and one based on time credits), relating their
soundness and relative-completeness proofs; our kernel sits closest to their
big-step reading, but carries a parametric cost model and a machine-checked
completeness proof.
The immediate ancestor of Silva's logic is the EasyCrypt cost
logic~\cite{Barbosa2021}, whose triples bound the running time of adversaries so
that concrete-security reductions can be discharged in the same prover as the
functional proof.
Relational versions of such logics, which bound the cost \emph{difference} between
two runs rather than a single run, have also been
developed~\cite{Cicek2017}, a direction orthogonal to the single-program bounds we
treat.
What our mechanization adds to this line is a fully machine-checked metatheory for
Silva's logic, including a relative-completeness result the source text leaves
open.

\paragraph{Type-based and automatic resource analysis.}
A complementary line replaces proof obligations with a type discipline that
\emph{infers} bounds.
Crary and Weirich~\cite{CraryWeirich2000} certify resource bounds through a type
system with explicit resource witnesses, and Hofmann and
Jost~\cite{HofmannJost2003} introduce the amortized potential-as-type technique,
assigning potential to data, that underlies automatic amortized resource analysis.
Resource Aware ML~\cite{Hoffmann2012,HoffmannDasWeng2017} turns that technique into
an inference engine, reducing the search for polynomial (and later also
exponential and logarithmic) bounds to linear programming over potential
annotations, and TiML~\cite{Wang2017} attaches time annotations to ML types that an
SMT backend discharges.
A lighter-weight variant embeds the accounting in the host language: Danielsson's
thunk monad~\cite{Danielsson2008} tracks the amortized cost of purely functional
data structures inside Agda's type checker.
On the real-time side, the survey of Wilhelm et al.~\cite{Wilhelm2008} maps the
worst-case execution-time tool landscape that motivates verified timing guarantees
in the first place.
These approaches favour automation and inferred, often asymptotic, bounds; our work
sits at the opposite end, trading automation for interactive proofs of concrete
bounds backed by a machine-checked metatheory.

\paragraph{Mechanized cost verification.}
The potential method for amortized analysis is due to Tarjan~\cite{Tarjan1985}.
Nipkow~\cite{Nipkow2015} mechanizes the amortized complexity of a suite of
functional data structures in Isabelle/HOL, and Atkey~\cite{Atkey2010} embeds
amortized resource analysis in separation logic through \emph{time credits}, a
resource that each computation step consumes.
Chargu\'eraud and Pottier~\cite{ChargueraudPottier2019} scale time credits in CFML
to a machine-checked \emph{inverse-Ackermann} bound for union-find, and M\'evel et
al.~\cite{MevelJourdanPottier2019} add time \emph{receipts} to Iris, dualising
credits so that lower bounds, and not only upper bounds, become provable.
Beyond data structures, concrete algorithmic results have been fully verified: van
der Weegen and McKinna~\cite{vanderWeegenMcKinna2008} machine-check the average-case
$\BigO(n\log n)$ bound of quicksort in Rocq.
These developments verify the cost of individual programs or structures; our focus
is instead the metatheory of a cost logic, the soundness and relative completeness
of the rules on top of which such examples are discharged.

\paragraph{Asymptotic complexity in proof assistants.}
Gu\'eneau et al.~\cite{Gueneau2018} make $\BigO$ a first-class object in Rocq,
modelling it as a domination relation between functions along a filter so that
multivariate asymptotics and the composition of $\BigO$ facts across a program
stay sound, on top of CFML time credits.
Zhan and Haslbeck~\cite{ZhanHaslbeck2018} verify asymptotic time complexity of
imperative programs in Isabelle by combining a refinement framework with an
Akra--Bazzi library, reaching divide-and-conquer algorithms such as merge sort and
Karatsuba multiplication.
McCarthy et al.~\cite{McCarthy2018} instead thread running time through a monad in
Rocq, so that a program's cost becomes part of its type and accumulates as it runs.
These works build general asymptotic infrastructure and treat $\BigO$ as the
primary specification.
Our asymptotic layer is deliberately lightweight and plays the opposite role: the
underlying logic proves concrete bounds, and the $\BigO$ assertions are thin,
derivable consequences that recover the familiar vocabulary without weakening what
was proved, which is why elementary $\Nat$-only definitions of the classes suffice
for the examples.

\paragraph{Mechanization in \Lean{} and relative completeness.}
\Lean{}~\cite{Moura2021} serves as both a dependently typed functional language and
an interactive prover, and its library Mathlib~\cite{Mathlib2020} supplies the
arithmetic, order theory, and, for the logarithmic example, the $\code{Nat.log}$
lemmas the cost proofs rely on.
The shallow embedding of assertions we adopt is a standard mechanization technique
that trades a fixed syntactic assertion language for the direct expressiveness of
the metalogic.
This is precisely what lets our completeness argument dispense with the
expressiveness hypothesis that Cook's classic relative-completeness
result~\cite{Cook1978} must assume: the strongest postcondition and the exact
run-cost of a loop are already \Lean{} predicates, so no arithmetization is needed
to name them (Section~\ref{sec:logic}).

\section{Conclusion}\label{sec:conclusion}

We have presented a \Lean{} mechanization of the upper-bound fragment of Silva's
resource-aware Hoare logic (a cost-instrumented semantics, the cost logic with
soundness, a new relative-completeness result, and its verification condition
generator), together with the correction of an unsoundness uncovered during
mechanization.
On top of this base we built an elementary theory of asymptotic cost assertions
and lifted it to the level of triples, so that a program can be certified to run
in $\BigO(g)$ time as a derivable weakening of a concrete bound.
Two examples exercise the framework: a division algorithm in $\BigO(n)$ and the
halving loop of binary search in $\BigO(\log n)$, the latter obtained from a
variant that advances by one against a geometrically shrinking value.

As future work, we plan to: (a)~bring the amortized and exact-cost fragments of
the dissertation, which the artifact already mechanizes, into the same uniform
account; (b)~complete the VCG completeness result
(Theorem~2 of~\cite{Silva2022}), which requires synthesizing loop annotations
from valid triples; and (c)~extend the asymptotic layer to multivariate and
super-linear classes (e.g.\ $\BigO(n^2)$) and connect it to
Mathlib's \code{Asymptotics.IsBigO}.

\subsubsection{Artifact Availability.}
The complete \Lean{} source of the formalization, including the asymptotic
extension and both examples, is available in a public repository~\cite{repo-logic}.
Beyond the upper-bound logic presented here, it also mechanizes the
amortized and exact-cost fragments of the dissertation, with their verification
condition generators.
The development builds against Lean~4.32.0 and Mathlib v4.32.0, and every
theorem depends only on the axioms \code{propext}, \code{Classical.choice}, and
\code{Quot.sound}.

\bibliographystyle{splncs04}
\bibliography{references}

\appendix

\section{Proof Sketches}\label{sec:appendix}

This appendix expands the metatheoretic results of
Sections~\ref{sec:background}--\ref{sec:examples} into proof sketches that follow
the mechanized \Lean{} arguments. Each sketch names the tactics or lemmas that
carry the load, so a reader can locate the corresponding step in the artifact.

\paragraph{Lemma~\ref{lem:det} (determinism).}
\emph{If $\evals{S}{\State}{t_1}{\State_1}$ and
$\evals{S}{\State}{t_2}{\State_2}$, then $t_1 = t_2$ and
$\State_1 = \State_2$.}

\begin{proof}
The proof is by induction on the first derivation $\evals{S}{\State}{t_1}{\State_1}$,
generalizing the cost $t_2$ and state $\State_2$ of the second, which is then
inverted by case analysis. When both derivations use the same rule, the induction
hypotheses equate the costs and states of the sub-runs and reflexivity closes the
goal; when they use different rules for the same command, the guard value
disagrees, since $\bval{b}\State$ cannot be simultaneously \code{true} and
\code{false}, so the case is vacuous. Only the two \code{if} and two \code{while}
constructors produce such mixed cases, all discharged by \code{simp\_all}. The cost
and state projections (\code{cost\_unique}, \code{state\_unique}) are immediate
corollaries.\qed
\end{proof}

\paragraph{Theorem~\ref{thm:sound} (soundness).}
\emph{If $\vdash \triple{P}{S}{Q}{T}$ then $\Valid\;\triple{P}{S}{Q}{T}$.}

\begin{proof}
By induction on the derivation of $\vdash \triple{P}{S}{Q}{T}$, unfolding validity
to produce, for each initial state, a terminating run whose cost meets the bound.
The atomic rules \textsc{Skip}, \textsc{Assign}, and \textsc{Array-Assign} each
exhibit the matching \code{Eval} constructor, with the charged cost equal to the
run's cost. \textsc{Seq} composes the two runs from the induction hypotheses; the
first sub-triple carries the extra conjunct $T_2 \le n$ into its postcondition,
which lets \code{omega} verify $t_1 + t_2 \le T_1\,\State + n$. \textsc{If}
branches on $\bval{b}\State$; the taken branch supplies the run, and
$t \le \max(T_1,T_2)$ together with the guard cost closes the bound. The
\code{while} case is the crux: after fixing the entry state and abbreviating
$k = f\,\State$, the proof proceeds by \emph{strong induction} on $k$. If the guard
is false, the exit premise gives $\tbcost{b}\State \le E\,\State$ and the
\textsc{While-F} run; if it is true, the body premise yields one iteration that
strictly decreases the variant ($f\,\State_1 < k$), preserves the invariant, and
respects the budget $E\,\State_1 + c \le n$, after which the strong induction
hypothesis handles the shorter remaining loop and \code{omega} sums the three cost
contributions within $E$. \textsc{Consequence} and the existential-precondition
rule are routine.\qed
\end{proof}

\paragraph{Theorem~\ref{thm:complete} (relative completeness).}
\emph{If $\Valid\;\triple{P}{S}{Q}{T}$ then $\vdash \triple{P}{S}{Q}{T}$.}

\begin{proof}
The argument, outlined after the theorem statement in Section~\ref{sec:logic},
runs as follows in the artifact. Using classical choice, the exact resource of a
run and the number of loop iterations are reified as functions \code{runCost} and
\code{steps}, well-defined because determinism (Lemma~\ref{lem:det}) and the
auxiliary \code{WhileSteps.unique} make the underlying data unique. The core is a
pointwise lemma \code{complete\_singleton}: given a concrete run
$\evals{S}{\State_0}{t}{\State'}$ with $Q\,\State'$, the triple
$\triple{\State = \State_0}{S}{Q}{\lambda\_.\,t}$ is derivable, by induction on
$S$. The atomic and structural cases mirror the semantics; the \code{while} case
instantiates the kernel rule with the strongest-postcondition invariant
$I \triangleq \lambda\State.\ \exists\,tt\,\State_f.\
\evals{\code{while}\;b\;\code{do}\;S}{\State}{tt}{\State_f} \wedge Q\,\State_f$,
variant $f \triangleq \code{steps}$, and remaining cost $E \triangleq \code{runCost}$,
all of which are already \Lean{} predicates and functions, so no assertion needs to
be arithmetized. This is exactly where the shallow embedding removes Cook's
expressiveness hypothesis. Finally, \code{complete} takes a valid triple, extracts
for each $\State_0 \vDash P$ its unique run (validity together with determinism),
applies \code{complete\_singleton}, weakens the exact cost to $T$, and reassembles
the pointwise family through the existential-precondition rule
(\code{of\_pointwise}). Choice is used precisely twice: in defining
\code{runCost}/\code{steps}, and in the classical \code{by\_cases} on
$P\,\State_0$ inside \code{of\_pointwise}.\qed
\end{proof}

\paragraph{Theorem~\ref{thm:vcg} (VCG soundness).}
\emph{If $c$ is $\sum$-free and every condition in $\code{VC}(c, Q)$ holds, then
the triple $\triple{\code{wpc}(c,Q)_1}{S}{Q}{\code{wpc}(c,Q)_2}$, where $S$ is the
program underlying $c$, is derivable and hence valid.}

\begin{proof}
The theorem factors through a key lemma \code{wpc\_sound}: if $c$ is $\sum$-free
and $\code{VC}(c,Q)$ holds, then
$\vdash \triple{\code{wpc}(c,Q)_1}{\code{strip}(c)}{Q}{\code{wpc}(c,Q)_2}$,
proved by induction on the $\sum$-freeness derivation of $c$. The atomic cases
apply the matching kernel rule. The \code{seq} case needs the second component's
cost to be state-independent, so that the compositional \textsc{Seq} rule can take
a constant $n$; this is the separate lemma \code{wpc\_cost\_const}, itself proved by
induction from the constancy of $\tacost{\cdot}$ and $\tbcost{\cdot}$ on
$\sum$-free expressions (\code{tacost\_const}, \code{tbcost\_const}), which is the
sole place the $\sum$-freeness hypothesis is consumed. The \code{if} case applies
the kernel \textsc{If} rule, its two branch preconditions coming from the
guard-indexed conjuncts of \code{wpc}. The \code{while} case instantiates the
derived thesis-form loop rule with the annotation $(I,f,N,t,\cst{b})$: the body
verification condition provides the per-iteration triple together with the cost
bound ${}\le t(k)$, and the variant-bound and guard-cost conditions supply the
remaining premises, so the charged cost
$\sum_{i<N} t(i) + (N+1)\cdot\cst{b}$ is met exactly. \code{vcg\_sound} then chains
\code{wpc\_sound} with the two outer components of $\code{VCG}$, namely the
entailment $P \Rightarrow \code{wpc}(c,Q)_1$ and the cost weakening
$\code{wpc}(c,Q)_2 \le T$, through the consequence rule.\qed
\end{proof}

\paragraph{The division example (\code{division\_valid}, \code{division\_bigO}).}
\emph{For every initial dividend $n$,
$\Valid\;\triple{r = n \wedge q = 0 \wedge m > 0 \wedge n \ge 0}
{\code{while}\;m \le r\;\code{do}\;(r := r - m;\ q := q + 1)}
{n = r + m\cdot q \wedge r < m}{\lambda\_.\,20n + 5}$;
and, as a corollary, $\Valid\;\triple{P}{S}{Q}{\BigO(\lambda n.\,n)}$.}

\begin{proof}
The concrete triple is obtained by applying \code{vcg\_sound}, through its validity
corollary, to a discharged \code{VCG} obligation, so the work lies in checking the
finitely many side conditions. With invariant
$I \equiv n = q\cdot m + r \wedge m>0 \wedge r\ge 0$, variant $f \equiv n - r$, and
bound $N \equiv n$, the body condition unfolds \code{wpc} and the substitution and
reduces to integer arithmetic closed by \code{omega} (using
$(q{+}1)\cdot m = q\cdot m + m$); the computed cost
$\sum_{i<n} 10 + (n{+}1)\cdot 3 = 13n + 3$, evaluated with \code{Finset.sum\_const}
and \code{Finset.card\_range}, is weakened to the stated $20n+5$. The asymptotic
corollary instantiates \code{BigOValid} with the bound read as a function of the
input size and transports the concrete validity across the $\Int$/$\Nat$ cast,
using $\lambda n.\,20n+5 \in \BigO(n)$ (\code{isBigO\_linear}).\qed
\end{proof}

\paragraph{The halving example (\code{halving\_valid}, \code{halving\_bigO}).}
\emph{For every initial value $n$,
$\Valid\;\triple{m = n \wedge n \ge 0}{\code{while}\;m > 1\;\code{do}\;m := m/2}
{m \le 1}{\lambda\_.\,7\log_2 n + 5}$;
and, as a corollary, $\Valid\;\triple{P}{S}{Q}{\BigO(\lambda n.\,\log_2 n)}$.}

\begin{proof}
As in the division example, the concrete triple follows from \code{vcg\_sound}
applied to a discharged \code{VCG} obligation. With invariant
$I \equiv 0 \le m \le n$, variant $f \equiv \log_2 n - \log_2 m$, and bound
$N \equiv \log_2 n$, the decisive body condition is that $f$ strictly increases; it
does, by exactly one, because $\log_2(m/2) = \log_2 m - 1$
(\code{Nat.log\_div\_base}), and \code{omega} closes the surrounding inequalities
from the monotonicity and positivity facts about \code{Nat.log}; the computed cost
$\sum_{i<\log_2 n} 4 + (\log_2 n {+} 1)\cdot 3 = 7\log_2 n + 3$ is weakened to
$7\log_2 n + 5$. The asymptotic corollary lifts the bound as before, now via
$\lambda n.\,7\log_2 n + 5 \in \BigO(\log n)$ (\code{isBigO\_log}).\qed
\end{proof}

\end{document}
